\section{A conditional route from iterator uniqueness to induction}
\label{sec:uniqueness-research}

This section records an arithmetic research direction, not a completed
replacement for the delivered checker.  Its positive result is a precise
retraction argument: scalar iterator uniqueness implies induction for
every predicate for which a suitable numerical conditional is already
available.  In particular, it implies induction for equations of the
form $t(n)=0$ using only numerical coverage and the ordinary zero test.
The missing step is an independent equality-reflection bootstrap for
the scalar language.  No new bi-interpretation or smaller complete
checker is claimed for the proposed uniqueness rule.

\subsection{The rule and its elementary consequences}

Uniqueness of recursion as an alternative arithmetic principle is
classical: Tait states the equivalence for a richer language with
counter-aware primitive recursion and product types \cite{tait}.
That result alone does not settle the restricted scalar update
$h(u)$ considered here. No novelty is claimed for the general idea
of using recursion uniqueness.

Work, for this arithmetic analysis, in ordinary intuitionistic
first-order logic with actual equality.  The numerical terms are
\[
 t ::= x\mid 0\mid S(t)\mid I(n,a;u.h),
\]
with capture-free substitution; $u$ is bound only in $h$.  Retain the
two computation equations
\begin{align*}
 I(0,a;u.h)&=a,\\
 I(Sn,a;u.h)&=h[I(n,a;u.h)/u].
\end{align*}
Terms such as $f(z)$ and $h(u)$ below are fixed finite term templates,
not variables of a function sort.  Passive parameters are suppressed.

Consider replacing literal-equation induction by the rule
\[
 \frac{\Gamma\vdash f(Sz)=h(f(z))}
      {\Gamma\vdash f(t)=I(t,f(0);u.h(u))}.
 \tag{UR}\label{rule:scalar-ur}
\]
Here $z$ is fresh for $\Gamma$, the target term $t$, and all passive
parameters; the update template $h(u)$ does not contain $z$.  Its only
active argument is the previous state $u$.  The displayed premise is
proved in $\Gamma$ itself: there is no discharged induction hypothesis.
The seed $f(0)$ is determined by $f$, so no separate base premise or
seed certificate occurs in this rule.

\paragraph{Equation induction proves UR.}
Assume the premise of UR and induct on the literal equation
\[
 f(n)=I(n,f(0);u.h(u)).
\]
The base is the zero computation equation and symmetry.  At a successor,
the recurrence premise, the current induction equation, congruence,
and the successor computation equation give
\begin{align*}
 f(Sn)&=h(f(n))\\
      &=h(I(n,f(0);u.h(u)))\\
      &=I(Sn,f(0);u.h(u)).
\end{align*}
Thus UR is sound in the old arithmetic.  This direction needs neither
coverage nor a numerical equality test.

\paragraph{Uniqueness with a common initial value.}
UR implies the following useful derived rule.  If, with fresh $z$,
\[
 \Gamma\vdash f(0)=g(0),\qquad
 \Gamma\vdash f(Sz)=h(f(z)),\qquad
 \Gamma\vdash g(Sz)=h(g(z)),
\]
then $\Gamma\vdash f(t)=g(t)$.  Apply UR to $f$ and $g$ separately
and use equality congruence in the seed of the iterator.  The argument
is contextual: arbitrary retained first-order assumptions may occur
in $\Gamma$, provided they do not contain the eigenvariable.

\paragraph{The successor iterator is the identity.}
Taking $f(z)=z$ and $h(u)=Su$, the premise of UR is reflexivity.
Consequently
\[
 \Gamma\vdash I(t,0;u.Su)=t.
 \tag{Id}\label{eq:ur-identity}
\]
This identity is a consequence of UR, not an additional axiom.  In
particular, the retraction argument below does not silently assume
that a successor orbit already exhausts the numerical domain.

\subsection{Retractions closed under successor are the identity}

\paragraph{Retraction lemma.}
Let $e(x)$ be a fixed scalar term template.  Suppose that $\Gamma$
proves the following equations, with $x$ fresh:
\begin{align}
 e(0)&=0,\label{eq:ur-e0}\\
 e(e(x))&=e(x),\label{eq:ur-idempotence}\\
 e(S(e(x)))&=S(e(x)).\label{eq:ur-closed}
\end{align}
Then the computation equations, equality reasoning, and UR prove
$\Gamma\vdash e(t)=t$ for every target term $t$.

\paragraph{Proof.}
Define two finite term abbreviations
\[
 h(u)=e(S(e(u))),\qquad r(n)=I(n,0;u.h(u)).
\]
No new operation or function quantifier is introduced.  Idempotence
implies the two global commutation equations
\begin{align}
 e(h(u))&=h(u),\label{eq:ur-eh}\\
 h(e(u))&=h(u).\label{eq:ur-he}
\end{align}
For the first, substitute $S(e(u))$ for $x$ in
\eqref{eq:ur-idempotence}.  For the second, unfold $h$ and replace
$e(e(u))$ by $e(u)$ under $S$ and $e$.  The successor-closure
assumption is exactly
\[
 h(u)=S(e(u)).
 \tag{H}\label{eq:ur-h-successor}
\]

Now consider $f(n)=e(r(n))$.  Its base is
$f(0)=e(0)=0$.  Its successor equation is, using the iterator
computation and \eqref{eq:ur-eh}--\eqref{eq:ur-he},
\[
 f(Sn)=e(h(r(n)))=h(r(n))=h(e(r(n)))=h(f(n)).
\]
UR, followed by replacement of the equal seeds, therefore yields
\[
 e(r(n))=I(n,0;u.h(u))=r(n).
 \tag{R1}\label{eq:ur-er}
\]
This is the first substantive use of uniqueness; it establishes that
the entire $h$ orbit lies in the image of the retraction.

Equation \eqref{eq:ur-h-successor} and \eqref{eq:ur-er} now give
\[
 r(Sn)=h(r(n))=S(e(r(n)))=S(r(n)),\qquad r(0)=0.
\]
Apply UR again, this time to $r$ with update $u\mapsto Su$.
By \eqref{eq:ur-identity},
\[
 r(n)=I(n,0;u.Su)=n.
 \tag{R2}\label{eq:ur-r-identity}
\]
Combining \eqref{eq:ur-er} and \eqref{eq:ur-r-identity} gives
$e(n)=n$.  Every displayed calculation is an equality derivation
under the fixed context.  No induction is used in the proof.
\hfill$\square$

\subsection{Induction from term-definable conditionals}

The preceding lemma gives a useful sufficient condition for induction.
Let $E(n)$ be a fixed formula.  Suppose that the retained context
proves its base and that its step is derivable:
\[
 \Gamma\vdash E(0),\qquad
 \Gamma,E(z)\vdash E(Sz),
 \tag{P}\label{eq:ur-ind-premises}
\]
where $z$ is fresh.  Suppose also that $E$ is decidable in this
context and has a term-defined split: there is a fixed term $e(n)$
such that, uniformly in fresh $n$,
\begin{align}
 \Gamma&\vdash E(n)\lor\neg E(n),\label{eq:ur-dec}\\
 \Gamma,E(n)&\vdash e(n)=n,\label{eq:ur-true-branch}\\
 \Gamma,\neg E(n)&\vdash e(n)=0.\label{eq:ur-false-branch}
\end{align}
Then UR proves the induction conclusion $\Gamma\vdash E(t)$.

Here is the complete reduction.  A case split using
\eqref{eq:ur-dec} proves $E(e(n))$: in the $E(n)$ branch use
\eqref{eq:ur-true-branch}; in the other branch use
\eqref{eq:ur-false-branch} and the base $E(0)$.  The true-branch
identity, instantiated at $e(n)$, gives
\[
 e(e(n))=e(n).
\]
The step from \eqref{eq:ur-ind-premises}, instantiated at $e(n)$,
gives $E(S(e(n)))$.  Another use of the true-branch identity yields
\[
 e(S(e(n)))=S(e(n)).
\]
Finally, the base and the true-branch identity at zero yield
$e(0)=0$.  Thus all three premises of the retraction lemma hold.
It follows that $e(n)=n$, and transporting the already proved
$E(e(n))$ along this equality gives $E(n)$, then $E(t)$.

This result applies only to predicates whose displayed numerical
split has been justified.  It does not supply induction for arbitrary
quantified formulas.  In particular, no general principle of deciding
first-order formulas or converting them to numerical tests is being
assumed.

\subsection{An unconditional application to zero equations}

For this subsection also assume numerical coverage and separation:
\[
 \forall x\,(x=0\lor\exists y\,x=Sy),\qquad 0=S0\to\bot.
\]
Define the selector, as in the inherited scalar language, by
\[
 C(v,a,b)=I(v,a;u.b).
\]
Its two endpoint equations are immediate computation:
\[
 C(0,a,b)=a,\qquad C(Sv,a,b)=b.
\]
Coverage supplies the corresponding case distinction for every
argument.  Moreover, $Sx=0$ is contradictory: applying the term
$v\mapsto C(v,0,S0)$ gives $S0=0$, and separation then gives
$\bot$.  Hence every zero equation $q(n)=0$ is decidable, without
constructing a test for equality of two arbitrary terms.

Take $E(n)$ to be the literal equation $q(n)=0$ and set
\[
 e(n)=C(q(n),n,0).
\]
If $E(n)$ holds, the zero selector equation gives $e(n)=n$.
If $\neg E(n)$ holds, coverage for $q(n)$ eliminates its zero
alternative; in its successor alternative the other selector
equation gives $e(n)=0$.  Thus
\eqref{eq:ur-dec}--\eqref{eq:ur-false-branch} have direct proofs.
The conditional induction result proves the following derived rule:
\[
 \frac{\Gamma\vdash q(0)=0\qquad
       \Gamma,q(z)=0\vdash q(Sz)=0}
      {\Gamma\vdash q(t)=0}.
 \tag{ZI}\label{rule:ur-zero-ind}
\]
Only coverage, separation, selector computation, equality reasoning,
and UR were needed.  No predecessor, pairing, multiplication,
general primitive-recursion compiler, or faithful equality
characteristic was used.

The distinction between zero equations and arbitrary equations is
essential.  Coverage decides whether a single number is zero.  It
does not itself give a term $\delta(a,b)$ satisfying
$\delta(a,b)=0\leftrightarrow a=b$ for arbitrary $a,b$.

\subsection{The remaining bootstrap and the comparison boundary}

If the UR-based arithmetic could independently construct a scalar
term $\delta(a,b)$ and prove
\[
 \delta(a,b)=0\quad\longleftrightarrow\quad a=b
 \tag{EqRef}\label{eq:ur-missing-reflection}
\]
for actual equality throughout its domain, then ZI would recover
literal-equation induction.  Given an induction on $s(n)=t(n)$,
apply ZI to $\delta(s(n),t(n))=0$ and transport its base, step,
and conclusion through \eqref{eq:ur-missing-reflection}.
Together with the earlier proof that equation induction implies UR,
this would establish equivalence of the two arithmetic principles
over the stated common background.

The existing scalar proof of \eqref{eq:ur-missing-reflection}
cannot simply be imported at this point.  Its predecessor and
reflection developments use literal-equation inductions with
nonconstant right sides, including reconstruction invariants.
Using those results to justify a replacement for that very induction
would be circular unless their proofs were translated to UR or ZI.
The argument above removes part of the obstacle by deriving ZI, but
does not complete that translation.

There is a separate logical bookkeeping issue for a proposed native
$Q,D$ kernel.  The proofs in this section use ordinary intuitionistic
equality logic, contextual case distinction, substitution, and cut.
When replacing a primitive rule in the native kernel, the required
native versions of these resources must be derived without appealing
to the removed rule.  The ordinary arithmetic lemma alone is not a
certificate of equivalence between the complete native calculi.

Finally, a premise with no local induction hypothesis need not produce
a smaller implementation.  UR replaces the hypothesis extension by
construction of $f(t)$, $f(0)$, and an iterator with update $h$.
Its complete syntax checks, eigenvariable checks, logical rules, and
public theorem interface would all have to be measured.  No such
complete smaller checker or completed bi-interpretation is asserted
here.  The precise further research problem is to finish the
equality-reflection bootstrap under UR, then perform a matched
implementation comparison of the resulting complete calculus.
